在 \(\Z\) 上考虑模 3 同余：两个整数如果除以 3 的余数相同，就归为一类。比如 \(1,4,7,10,\ldots\) 余数都是 1，它们构成一个等价类；\(0,3,6,9,\ldots\) 构成另一个等价类。总共三个等价类，恰好把全体整数分成三个互不相交的区块——这就是等价关系做分类的基本画面。

逻辑告诉我们怎样表述判断，集合则回答``正在谈论哪些对象''。从两个集合的笛卡尔积中挑出若干有序对，就得到关系；若每个输入恰好对应一个输出，这个关系又成为函数。三者逐层累加结构。

阅读顺序如下：

\begin{enumerate}
\tightlist
\item
  集合与集合运算：元素、子集、运算、幂集与笛卡尔积。
\item
  关系与等价类：用有序对描述``相等到某种程度''，理解等价类如何形成划分。
\item
  偏序与可比较性：区分先后、包含与大小，理解不必所有对象都能比较。
\item
  函数与可逆性：从特殊关系出发理解定义域、值域、单射、满射、复合与逆。
\end{enumerate}

这些概念会反复出现在后续课程中：等价关系通过商集构造新对象，偏序支撑上确界与递归结构，函数是微积分和概率论的基本对象。

\subsection{怎么读这一章}

按``对象--对象间关系--顺序--映射''阅读。集合说明有哪些对象，等价关系负责归类，偏序只比较一部分对象，函数则追踪输入怎样变成输出。每一篇都优先看反例，避免把相近概念误当成同一个。

\subsection{本章篇目}

以下篇目按概念依赖排列；若从具体问题进入，也可以先打开最接近当前困惑的一篇，再沿篇内链接回补。

\begin{itemize}
\tightlist
\item \hyperref[sec:01-Mathematical-Language/02-Sets-Relations-Functions/01]{``集合与集合运算''}；
\item \hyperref[sec:01-Mathematical-Language/02-Sets-Relations-Functions/02]{``关系与等价类''}；
\item \hyperref[sec:01-Mathematical-Language/02-Sets-Relations-Functions/03]{``偏序与可比较性''}；
\item \hyperref[sec:01-Mathematical-Language/02-Sets-Relations-Functions/04]{``函数与可逆性''}。
\end{itemize}
